%% Generated by Sphinx.
\def\sphinxdocclass{report}
\documentclass[letterpaper,10pt,english]{sphinxmanual}
\ifdefined\pdfpxdimen
   \let\sphinxpxdimen\pdfpxdimen\else\newdimen\sphinxpxdimen
\fi \sphinxpxdimen=.75bp\relax
\ifdefined\pdfimageresolution
    \pdfimageresolution= \numexpr \dimexpr1in\relax/\sphinxpxdimen\relax
\fi
%% let collapsible pdf bookmarks panel have high depth per default
\PassOptionsToPackage{bookmarksdepth=5}{hyperref}

\PassOptionsToPackage{booktabs}{sphinx}
\PassOptionsToPackage{colorrows}{sphinx}

\PassOptionsToPackage{warn}{textcomp}
\usepackage[utf8]{inputenc}
\ifdefined\DeclareUnicodeCharacter
% support both utf8 and utf8x syntaxes
  \ifdefined\DeclareUnicodeCharacterAsOptional
    \def\sphinxDUC#1{\DeclareUnicodeCharacter{"#1}}
  \else
    \let\sphinxDUC\DeclareUnicodeCharacter
  \fi
  \sphinxDUC{00A0}{\nobreakspace}
  \sphinxDUC{2500}{\sphinxunichar{2500}}
  \sphinxDUC{2502}{\sphinxunichar{2502}}
  \sphinxDUC{2514}{\sphinxunichar{2514}}
  \sphinxDUC{251C}{\sphinxunichar{251C}}
  \sphinxDUC{2572}{\textbackslash}
\fi
\usepackage{cmap}
\usepackage[T1]{fontenc}
\usepackage{amsmath,amssymb,amstext}
\usepackage{babel}



\usepackage{tgtermes}
\usepackage{tgheros}
\renewcommand{\ttdefault}{txtt}



\usepackage[Bjarne]{fncychap}
\usepackage{sphinx}

\fvset{fontsize=auto}
\usepackage{geometry}


% Include hyperref last.
\usepackage{hyperref}
% Fix anchor placement for figures with captions.
\usepackage{hypcap}% it must be loaded after hyperref.
% Set up styles of URL: it should be placed after hyperref.
\urlstyle{same}

\addto\captionsenglish{\renewcommand{\contentsname}{Contents:}}

\usepackage{sphinxmessages}
\setcounter{tocdepth}{1}



\title{hammindist}
\date{May 25, 2026}
\release{0.1.0}
\author{Students}
\newcommand{\sphinxlogo}{\vbox{}}
\renewcommand{\releasename}{Release}
\makeindex
\begin{document}

\ifdefined\shorthandoff
  \ifnum\catcode`\=\string=\active\shorthandoff{=}\fi
  \ifnum\catcode`\"=\active\shorthandoff{"}\fi
\fi

\pagestyle{empty}
\sphinxmaketitle
\pagestyle{plain}
\sphinxtableofcontents
\pagestyle{normal}
\phantomsection\label{\detokenize{index::doc}}



\chapter{Solvers API}
\label{\detokenize{index:module-hammindist.functions}}\label{\detokenize{index:solvers-api}}\index{module@\spxentry{module}!hammindist.functions@\spxentry{hammindist.functions}}\index{hammindist.functions@\spxentry{hammindist.functions}!module@\spxentry{module}}\index{A() (in module hammindist.functions)@\spxentry{A()}\spxextra{in module hammindist.functions}}

\begin{fulllineitems}
\phantomsection\label{\detokenize{index:hammindist.functions.A}}
\pysigstartsignatures
\pysiglinewithargsret
{\sphinxcode{\sphinxupquote{hammindist.functions.}}\sphinxbfcode{\sphinxupquote{A}}}
{\sphinxparam{\DUrole{n}{n}}\sphinxparamcomma \sphinxparam{\DUrole{n}{d}}\sphinxparamcomma \sphinxparam{\DUrole{n}{verbose}\DUrole{o}{=}\DUrole{default_value}{True}}}
{}
\pysigstopsignatures
\sphinxAtStartPar
Calculate A(n, d)
:param n: the length of the code
:type n: int
:param d: the minimum distance in the final set
:type d: int
\begin{quote}\begin{description}
\sphinxlineitem{Returns}
\sphinxAtStartPar
the length of the set (number of elements)
code (set): the set with the elements at minimum distance d

\sphinxlineitem{Return type}
\sphinxAtStartPar
size (int)

\end{description}\end{quote}

\end{fulllineitems}

\index{HammingTupla (class in hammindist.functions)@\spxentry{HammingTupla}\spxextra{class in hammindist.functions}}

\begin{fulllineitems}
\phantomsection\label{\detokenize{index:hammindist.functions.HammingTupla}}
\pysigstartsignatures
\pysiglinewithargsret
{\sphinxbfcode{\sphinxupquote{\DUrole{k}{class}\DUrole{w}{ }}}\sphinxcode{\sphinxupquote{hammindist.functions.}}\sphinxbfcode{\sphinxupquote{HammingTupla}}}
{\sphinxparam{\DUrole{n}{lenght}}\sphinxparamcomma \sphinxparam{\DUrole{n}{distance}}}
{}
\pysigstopsignatures
\sphinxAtStartPar
A class to generate binary tuples of a given length with Hamming weight
less than or equal to a specified distance.
\index{lenght (hammindist.functions.HammingTupla attribute)@\spxentry{lenght}\spxextra{hammindist.functions.HammingTupla attribute}}

\begin{fulllineitems}
\phantomsection\label{\detokenize{index:hammindist.functions.HammingTupla.lenght}}
\pysigstartsignatures
\pysiglinewithargsret
{\sphinxbfcode{\sphinxupquote{lenght}}}
{\sphinxparam{\DUrole{n}{int}}}
{}
\pysigstopsignatures
\sphinxAtStartPar
The lenght of the binary tuples.

\end{fulllineitems}

\index{distance (hammindist.functions.HammingTupla attribute)@\spxentry{distance}\spxextra{hammindist.functions.HammingTupla attribute}}

\begin{fulllineitems}
\phantomsection\label{\detokenize{index:hammindist.functions.HammingTupla.distance}}
\pysigstartsignatures
\pysiglinewithargsret
{\sphinxbfcode{\sphinxupquote{distance}}}
{\sphinxparam{\DUrole{n}{int}}}
{}
\pysigstopsignatures
\sphinxAtStartPar
The maximum number of 1s (Hamming weight) allowed in each tuple.

\end{fulllineitems}

\index{graph (hammindist.functions.HammingTupla attribute)@\spxentry{graph}\spxextra{hammindist.functions.HammingTupla attribute}}

\begin{fulllineitems}
\phantomsection\label{\detokenize{index:hammindist.functions.HammingTupla.graph}}
\pysigstartsignatures
\pysiglinewithargsret
{\sphinxbfcode{\sphinxupquote{graph}}}
{\sphinxparam{\DUrole{n}{networkx.Graph}}}
{}
\pysigstopsignatures
\sphinxAtStartPar
The graph representation of the Hamming tuples, where nodes are tuples
and edges connect tuples with Hamming distance 1.

\end{fulllineitems}

\index{as\_list() (hammindist.functions.HammingTupla method)@\spxentry{as\_list()}\spxextra{hammindist.functions.HammingTupla method}}

\begin{fulllineitems}
\phantomsection\label{\detokenize{index:hammindist.functions.HammingTupla.as_list}}
\pysigstartsignatures
\pysiglinewithargsret
{\sphinxbfcode{\sphinxupquote{as\_list}}}
{}
{}
\pysigstopsignatures
\sphinxAtStartPar
Return all generated tuples as a list.

\sphinxAtStartPar
Returns:
list of tuple of int
\begin{quote}

\sphinxAtStartPar
A list containing all generated binary tuples.
\end{quote}

\end{fulllineitems}

\index{draw() (hammindist.functions.HammingTupla method)@\spxentry{draw()}\spxextra{hammindist.functions.HammingTupla method}}

\begin{fulllineitems}
\phantomsection\label{\detokenize{index:hammindist.functions.HammingTupla.draw}}
\pysigstartsignatures
\pysiglinewithargsret
{\sphinxbfcode{\sphinxupquote{draw}}}
{}
{}
\pysigstopsignatures
\sphinxAtStartPar
Render the Hamming graph using Matplotlib without node labels.

\sphinxAtStartPar
Useful for larger graphs where text labels would clutter the
visualization. Nodes are drawn as small ‘orchid’ colored circles.

\end{fulllineitems}

\index{get\_instances\_complete() (hammindist.functions.HammingTupla method)@\spxentry{get\_instances\_complete()}\spxextra{hammindist.functions.HammingTupla method}}

\begin{fulllineitems}
\phantomsection\label{\detokenize{index:hammindist.functions.HammingTupla.get_instances_complete}}
\pysigstartsignatures
\pysiglinewithargsret
{\sphinxbfcode{\sphinxupquote{get\_instances\_complete}}}
{}
{}
\pysigstopsignatures
\sphinxAtStartPar
Generate all binary tuples of the given lenght (2\textasciicircum{}lenght tuples).
This method is independent of the distance attribute.
\begin{quote}\begin{description}
\sphinxlineitem{Yields}
\sphinxAtStartPar
\sphinxstyleemphasis{tuple of int} \textendash{} A binary tuple (containing 0s and 1s).

\end{description}\end{quote}

\end{fulllineitems}

\index{get\_instances\_restricted() (hammindist.functions.HammingTupla method)@\spxentry{get\_instances\_restricted()}\spxextra{hammindist.functions.HammingTupla method}}

\begin{fulllineitems}
\phantomsection\label{\detokenize{index:hammindist.functions.HammingTupla.get_instances_restricted}}
\pysigstartsignatures
\pysiglinewithargsret
{\sphinxbfcode{\sphinxupquote{get\_instances\_restricted}}}
{}
{}
\pysigstopsignatures
\end{fulllineitems}

\index{hamming\_distance() (hammindist.functions.HammingTupla static method)@\spxentry{hamming\_distance()}\spxextra{hammindist.functions.HammingTupla static method}}

\begin{fulllineitems}
\phantomsection\label{\detokenize{index:hammindist.functions.HammingTupla.hamming_distance}}
\pysigstartsignatures
\pysiglinewithargsret
{\sphinxbfcode{\sphinxupquote{\DUrole{k}{static}\DUrole{w}{ }}}\sphinxbfcode{\sphinxupquote{hamming\_distance}}}
{\sphinxparam{\DUrole{n}{X}}\sphinxparamcomma \sphinxparam{\DUrole{n}{Y}}}
{}
\pysigstopsignatures
\sphinxAtStartPar
Calculate the distance between two elements with the same lenght
Arg:
\begin{quote}

\sphinxAtStartPar
X (List{[}any{]}): a string, an array or a tuple of 0’s and 1’s with lenght n
Y (List{[}any{]}): a string, an array or a tuple of 0’s and 1’s with lenght n
\end{quote}
\begin{quote}\begin{description}
\sphinxlineitem{Returns}
\sphinxAtStartPar
a integer value non negative

\sphinxlineitem{Return type}
\sphinxAtStartPar
d (int)

\end{description}\end{quote}

\end{fulllineitems}


\end{fulllineitems}

\index{accept() (in module hammindist.functions)@\spxentry{accept()}\spxextra{in module hammindist.functions}}

\begin{fulllineitems}
\phantomsection\label{\detokenize{index:hammindist.functions.accept}}
\pysigstartsignatures
\pysiglinewithargsret
{\sphinxcode{\sphinxupquote{hammindist.functions.}}\sphinxbfcode{\sphinxupquote{accept}}}
{\sphinxparam{\DUrole{n}{current\_code}}\sphinxparamcomma \sphinxparam{\DUrole{n}{candidate}}\sphinxparamcomma \sphinxparam{\DUrole{n}{T}}}
{}
\pysigstopsignatures
\sphinxAtStartPar
Simulated annealing acceptance criterion.
\begin{itemize}
\item {} 
\sphinxAtStartPar
Improvements are always accepted.

\item {} 
\sphinxAtStartPar
Worse solutions may be accepted with probability
exp(Δ/T), where Δ is the size difference between solutions.

\end{itemize}
\begin{quote}\begin{description}
\sphinxlineitem{Parameters}\begin{itemize}
\item {} 
\sphinxAtStartPar
\sphinxstyleliteralstrong{\sphinxupquote{current\_code}} (\sphinxstyleliteralemphasis{\sphinxupquote{List}}\sphinxstyleliteralemphasis{\sphinxupquote{{[}}}\sphinxstyleliteralemphasis{\sphinxupquote{Any}}\sphinxstyleliteralemphasis{\sphinxupquote{{]}}}) \textendash{} Current solution (list of codewords).
Typically the best\sphinxhyphen{}known or current
working solution.

\item {} 
\sphinxAtStartPar
\sphinxstyleliteralstrong{\sphinxupquote{candidate}} (\sphinxstyleliteralemphasis{\sphinxupquote{List}}\sphinxstyleliteralemphasis{\sphinxupquote{{[}}}\sphinxstyleliteralemphasis{\sphinxupquote{Any}}\sphinxstyleliteralemphasis{\sphinxupquote{{]}}}) \textendash{} Candidate (neighbor) solution being evaluated.
Usually generated by a perturbation or
local move from the current solution.

\item {} 
\sphinxAtStartPar
\sphinxstyleliteralstrong{\sphinxupquote{T}} (\sphinxstyleliteralemphasis{\sphinxupquote{float}}) \textendash{} Current temperature parameter in simulated annealing.
Higher temperatures increase the probability of accepting
worse solutions. As T decreases toward 0, the algorithm
becomes more greedy (only accepts improvements).

\end{itemize}

\sphinxlineitem{Returns}
\sphinxAtStartPar
\begin{description}
\sphinxlineitem{Either the candidate solution (if accepted) or the}
\sphinxAtStartPar
current\_code (if rejected). Returns candidate for
improvements or when probabilistic acceptance succeeds.

\end{description}


\sphinxlineitem{Return type}
\sphinxAtStartPar
List{[}Any{]}

\end{description}\end{quote}

\end{fulllineitems}

\index{binary\_hamming\_distance() (in module hammindist.functions)@\spxentry{binary\_hamming\_distance()}\spxextra{in module hammindist.functions}}

\begin{fulllineitems}
\phantomsection\label{\detokenize{index:hammindist.functions.binary_hamming_distance}}
\pysigstartsignatures
\pysiglinewithargsret
{\sphinxcode{\sphinxupquote{hammindist.functions.}}\sphinxbfcode{\sphinxupquote{binary\_hamming\_distance}}}
{\sphinxparam{\DUrole{n}{x}\DUrole{p}{:}\DUrole{w}{ }\DUrole{n}{int}}\sphinxparamcomma \sphinxparam{\DUrole{n}{y}\DUrole{p}{:}\DUrole{w}{ }\DUrole{n}{int}}}
{{ $\rightarrow$ int}}
\pysigstopsignatures
\end{fulllineitems}

\index{build\_adjacency\_bitsets() (in module hammindist.functions)@\spxentry{build\_adjacency\_bitsets()}\spxextra{in module hammindist.functions}}

\begin{fulllineitems}
\phantomsection\label{\detokenize{index:hammindist.functions.build_adjacency_bitsets}}
\pysigstartsignatures
\pysiglinewithargsret
{\sphinxcode{\sphinxupquote{hammindist.functions.}}\sphinxbfcode{\sphinxupquote{build\_adjacency\_bitsets}}}
{\sphinxparam{\DUrole{n}{n}\DUrole{p}{:}\DUrole{w}{ }\DUrole{n}{int}}\sphinxparamcomma \sphinxparam{\DUrole{n}{d}\DUrole{p}{:}\DUrole{w}{ }\DUrole{n}{int}}\sphinxparamcomma \sphinxparam{\DUrole{n}{even\_only}\DUrole{p}{:}\DUrole{w}{ }\DUrole{n}{bool}\DUrole{w}{ }\DUrole{o}{=}\DUrole{w}{ }\DUrole{default_value}{False}}}
{}
\pysigstopsignatures
\sphinxAtStartPar
Build an adjacency bitset representation of a graph where vertices are binary strings
of length n, and edges connect vertices whose Hamming distance is at least d.
:param n: Length of binary strings (number of bits).
:type n: int
:param d: Minimum Hamming distance required for an edge between vertices.
:type d: int
:param even\_only: If True, only include vertices with even parity
\begin{quote}

\sphinxAtStartPar
(even number of 1\sphinxhyphen{}bits). Defaults to False.
\end{quote}
\begin{quote}\begin{description}
\sphinxlineitem{Returns}
\sphinxAtStartPar
\begin{description}
\sphinxlineitem{Adjacency list as bitsets. adj{[}i{]} is an integer whose}
\sphinxAtStartPar
j\sphinxhyphen{}th bit (from LSB) indicates if vertex i is connected to vertex j.
Only bits for j \textgreater{} i are stored (lower triangular representation).

\end{description}

\sphinxAtStartPar
N (int): Number of vertices in the graph.
vertices (List{[}int{]}): List of vertex values (integers) in the same order
\begin{quote}

\sphinxAtStartPar
as their indices in the adjacency list.
\end{quote}


\sphinxlineitem{Return type}
\sphinxAtStartPar
adj (List{[}int{]})

\end{description}\end{quote}

\end{fulllineitems}

\index{build\_outside\_pool() (in module hammindist.functions)@\spxentry{build\_outside\_pool()}\spxextra{in module hammindist.functions}}

\begin{fulllineitems}
\phantomsection\label{\detokenize{index:hammindist.functions.build_outside_pool}}
\pysigstartsignatures
\pysiglinewithargsret
{\sphinxcode{\sphinxupquote{hammindist.functions.}}\sphinxbfcode{\sphinxupquote{build\_outside\_pool}}}
{\sphinxparam{\DUrole{n}{all\_words\_set}}\sphinxparamcomma \sphinxparam{\DUrole{n}{candidate}}\sphinxparamcomma \sphinxparam{\DUrole{n}{tabu\_expiry}}\sphinxparamcomma \sphinxparam{\DUrole{n}{step}}}
{}
\pysigstopsignatures
\sphinxAtStartPar
Build a pool of candidate codewords for insertion.

\sphinxAtStartPar
Only codewords that:
\sphinxhyphen{} Are not part of the current candidate solution.
\sphinxhyphen{} Are not tabu.

\sphinxAtStartPar
The resulting pool is shuffled to encourage exploration
of the search space.
\begin{quote}\begin{description}
\sphinxlineitem{Parameters}\begin{itemize}
\item {} 
\sphinxAtStartPar
\sphinxstyleliteralstrong{\sphinxupquote{all\_words\_set}} (\sphinxstyleliteralemphasis{\sphinxupquote{Set}}\sphinxstyleliteralemphasis{\sphinxupquote{{[}}}\sphinxstyleliteralemphasis{\sphinxupquote{int}}\sphinxstyleliteralemphasis{\sphinxupquote{{]}}}) \textendash{} Set of all possible codewords in the
search space (e.g., all binary strings
of length n that satisfy constraints).

\item {} 
\sphinxAtStartPar
\sphinxstyleliteralstrong{\sphinxupquote{candidate}} (\sphinxstyleliteralemphasis{\sphinxupquote{List}}\sphinxstyleliteralemphasis{\sphinxupquote{{[}}}\sphinxstyleliteralemphasis{\sphinxupquote{int}}\sphinxstyleliteralemphasis{\sphinxupquote{{]}}}) \textendash{} Current candidate solution containing
codewords already selected.

\item {} 
\sphinxAtStartPar
\sphinxstyleliteralstrong{\sphinxupquote{tabu\_expiry}} (\sphinxstyleliteralemphasis{\sphinxupquote{Dict}}\sphinxstyleliteralemphasis{\sphinxupquote{{[}}}\sphinxstyleliteralemphasis{\sphinxupquote{int}}\sphinxstyleliteralemphasis{\sphinxupquote{, }}\sphinxstyleliteralemphasis{\sphinxupquote{int}}\sphinxstyleliteralemphasis{\sphinxupquote{{]}}}) \textendash{} Dictionary mapping codewords to the
step number when their tabu status
expires. A codeword is considered
tabu if its expiry step \textgreater{} current step.

\item {} 
\sphinxAtStartPar
\sphinxstyleliteralstrong{\sphinxupquote{step}} (\sphinxstyleliteralemphasis{\sphinxupquote{int}}) \textendash{} Current iteration step number. Used to determine
which codewords have expired tabu status.

\end{itemize}

\sphinxlineitem{Returns}
\sphinxAtStartPar
\begin{description}
\sphinxlineitem{A shuffled list of codewords eligible for insertion.}
\sphinxAtStartPar
Returns an empty list if no eligible codewords exist.

\end{description}


\sphinxlineitem{Return type}
\sphinxAtStartPar
List{[}int{]}

\end{description}\end{quote}

\end{fulllineitems}

\index{greedy\_color() (in module hammindist.functions)@\spxentry{greedy\_color()}\spxextra{in module hammindist.functions}}

\begin{fulllineitems}
\phantomsection\label{\detokenize{index:hammindist.functions.greedy_color}}
\pysigstartsignatures
\pysiglinewithargsret
{\sphinxcode{\sphinxupquote{hammindist.functions.}}\sphinxbfcode{\sphinxupquote{greedy\_color}}}
{\sphinxparam{\DUrole{n}{P\_bits}}\sphinxparamcomma \sphinxparam{\DUrole{n}{adj}}}
{}
\pysigstopsignatures
\sphinxAtStartPar
Greedy coloring of the subgraph induced by a set of vertices.

\sphinxAtStartPar
This function performs a greedy graph coloring on the subgraph containing
only the vertices specified by P\_bits. The coloring order is optimized
by sorting vertices by their degree within the induced subgraph (highest
degree first).
\begin{quote}\begin{description}
\sphinxlineitem{Parameters}\begin{itemize}
\item {} 
\sphinxAtStartPar
\sphinxstyleliteralstrong{\sphinxupquote{P\_bits}} (\sphinxstyleliteralemphasis{\sphinxupquote{int}}) \textendash{} Bitset representing the set of vertices to include in the
subgraph. Bit i is set if vertex i is present.

\item {} 
\sphinxAtStartPar
\sphinxstyleliteralstrong{\sphinxupquote{adj}} (\sphinxstyleliteralemphasis{\sphinxupquote{List}}\sphinxstyleliteralemphasis{\sphinxupquote{{[}}}\sphinxstyleliteralemphasis{\sphinxupquote{int}}\sphinxstyleliteralemphasis{\sphinxupquote{{]}}}) \textendash{} Adjacency bitsets for the full graph. adj{[}i{]} is an
integer where bit j indicates an edge between vertices
i and j (assuming j \textgreater{} i for upper triangular storage).

\end{itemize}

\sphinxlineitem{Returns}
\sphinxAtStartPar
\begin{description}
\sphinxlineitem{List of vertices sorted by color in}
\sphinxAtStartPar
descending order (vertices with higher color numbers first).

\end{description}

\sphinxAtStartPar
num\_colors (int): Number of colors used in the coloring.


\sphinxlineitem{Return type}
\sphinxAtStartPar
sorted\_verts (List{[}int{]})

\end{description}\end{quote}

\end{fulllineitems}

\index{greedy\_initial\_clique() (in module hammindist.functions)@\spxentry{greedy\_initial\_clique()}\spxextra{in module hammindist.functions}}

\begin{fulllineitems}
\phantomsection\label{\detokenize{index:hammindist.functions.greedy_initial_clique}}
\pysigstartsignatures
\pysiglinewithargsret
{\sphinxcode{\sphinxupquote{hammindist.functions.}}\sphinxbfcode{\sphinxupquote{greedy\_initial\_clique}}}
{\sphinxparam{\DUrole{n}{adj}}\sphinxparamcomma \sphinxparam{\DUrole{n}{N}}\sphinxparamcomma \sphinxparam{\DUrole{n}{vertices}}}
{}
\pysigstopsignatures
\sphinxAtStartPar
Find a large clique using a greedy heuristic to obtain a lower bound.
:param adj: Adjacency bitsets for the graph. adj{[}i{]} is an integer
\begin{quote}

\sphinxAtStartPar
where bit j indicates an edge between vertices i and j.
\end{quote}
\begin{quote}\begin{description}
\sphinxlineitem{Parameters}\begin{itemize}
\item {} 
\sphinxAtStartPar
\sphinxstyleliteralstrong{\sphinxupquote{N}} \textendash{} Number of vertices in the graph (length of adj list).

\item {} 
\sphinxAtStartPar
\sphinxstyleliteralstrong{\sphinxupquote{vertices}} \textendash{} List of vertex values (integers) corresponding
to indices 0..N\sphinxhyphen{}1 in the adjacency list. This
parameter is primarily for interface consistency
with other functions, but not used directly in
the heuristic.

\end{itemize}

\sphinxlineitem{Returns}
\sphinxAtStartPar

\sphinxAtStartPar
Size of the found clique (number of vertices).
clique\_bitset (int): Integer bitset where bit i is set if vertex i
\begin{quote}

\sphinxAtStartPar
(index in adjacency list) belongs to the clique.
\end{quote}


\sphinxlineitem{Return type}
\sphinxAtStartPar
size (int)

\end{description}\end{quote}

\end{fulllineitems}

\index{greedy\_start() (in module hammindist.functions)@\spxentry{greedy\_start()}\spxextra{in module hammindist.functions}}

\begin{fulllineitems}
\phantomsection\label{\detokenize{index:hammindist.functions.greedy_start}}
\pysigstartsignatures
\pysiglinewithargsret
{\sphinxcode{\sphinxupquote{hammindist.functions.}}\sphinxbfcode{\sphinxupquote{greedy\_start}}}
{\sphinxparam{\DUrole{n}{all\_words}}\sphinxparamcomma \sphinxparam{\DUrole{n}{d}}}
{}
\pysigstopsignatures
\sphinxAtStartPar
Build an initial valid code using a greedy random approach.

\sphinxAtStartPar
Shuffles all candidate codewords and adds each one to the code
if it satisfies the minimum distance constraint with all
already\sphinxhyphen{}selected codewords.

\sphinxAtStartPar
This is used as a warm start for more sophisticated heuristics
such as simulated annealing.

\sphinxAtStartPar
Args:
all\_words (List{[}int{]})
\begin{quote}

\sphinxAtStartPar
All candidate binary codewords of length n.
\end{quote}
\begin{description}
\sphinxlineitem{d (int):}
\sphinxAtStartPar
Minimum Hamming distance required between any two codewords.

\end{description}

\sphinxAtStartPar
Returns:
code (List{[}int{]}):
\begin{quote}

\sphinxAtStartPar
A valid code built greedily. Not guaranteed to be optimal.
\end{quote}

\end{fulllineitems}

\index{heuristic() (in module hammindist.functions)@\spxentry{heuristic()}\spxextra{in module hammindist.functions}}

\begin{fulllineitems}
\phantomsection\label{\detokenize{index:hammindist.functions.heuristic}}
\pysigstartsignatures
\pysiglinewithargsret
{\sphinxcode{\sphinxupquote{hammindist.functions.}}\sphinxbfcode{\sphinxupquote{heuristic}}}
{\sphinxparam{\DUrole{n}{n}}\sphinxparamcomma \sphinxparam{\DUrole{n}{d}}\sphinxparamcomma \sphinxparam{\DUrole{n}{iterations}\DUrole{o}{=}\DUrole{default_value}{1000}}\sphinxparamcomma \sphinxparam{\DUrole{n}{seed}\DUrole{o}{=}\DUrole{default_value}{None}}}
{}
\pysigstopsignatures
\sphinxAtStartPar
Estimate A(n, d) — the maximum size of a binary code of length n
and minimum Hamming distance d — using a greedy heuristic with
random restarts.

\sphinxAtStartPar
The algorithm builds a valid code greedily: at each step it tries
to add a random codeword that satisfies the minimum distance
constraint with all already\sphinxhyphen{}selected codewords. When no more
codewords can be added, it records the size and restarts.

\sphinxAtStartPar
This is a heuristic, so the result is a lower bound on A(n, d).
It is not guaranteed to find the true optimum.

\sphinxAtStartPar
Args:
n (int):
\begin{quote}

\sphinxAtStartPar
Length of the binary codewords.
\end{quote}
\begin{description}
\sphinxlineitem{d (int):}
\sphinxAtStartPar
Minimum Hamming distance required between any two codewords.

\sphinxlineitem{iterations}{[}int, optional{]}
\sphinxAtStartPar
Number of random restarts (default is 1000).

\sphinxlineitem{seed (int or None)}
\sphinxAtStartPar
Random seed for reproducibility (default is None).

\end{description}

\sphinxAtStartPar
Returns:
\begin{description}
\sphinxlineitem{best\_code (List{[}int{]}):}
\sphinxAtStartPar
The largest valid code found.

\end{description}

\end{fulllineitems}

\index{is\_valid\_set() (in module hammindist.functions)@\spxentry{is\_valid\_set()}\spxextra{in module hammindist.functions}}

\begin{fulllineitems}
\phantomsection\label{\detokenize{index:hammindist.functions.is_valid_set}}
\pysigstartsignatures
\pysiglinewithargsret
{\sphinxcode{\sphinxupquote{hammindist.functions.}}\sphinxbfcode{\sphinxupquote{is\_valid\_set}}}
{\sphinxparam{\DUrole{n}{candidate\_set}}\sphinxparamcomma \sphinxparam{\DUrole{n}{new\_code}}\sphinxparamcomma \sphinxparam{\DUrole{n}{min\_distance}}}
{}
\pysigstopsignatures
\sphinxAtStartPar
Valid if a candidate\_set is a valid set, it means that the elements of the set
are at least the required distance apart
:param candidate\_set: a set of codes with min\_distance between them
:type candidate\_set: set
:param new\_code: a code that will be added to the set if it has min\_distance with all
\begin{quote}

\sphinxAtStartPar
the elements of th candidate\_set
\end{quote}
\begin{quote}\begin{description}
\sphinxlineitem{Parameters}
\sphinxAtStartPar
\sphinxstyleliteralstrong{\sphinxupquote{min\_distance}} (\sphinxstyleliteralemphasis{\sphinxupquote{int}}) \textendash{} the minimum distance that must havee the all elements of the set
with the new code

\sphinxlineitem{Returns}
\sphinxAtStartPar
True if the set is valid, False if the set is not valid

\sphinxlineitem{Return type}
\sphinxAtStartPar
bool

\end{description}\end{quote}

\end{fulllineitems}

\index{max\_clique\_tomita() (in module hammindist.functions)@\spxentry{max\_clique\_tomita()}\spxextra{in module hammindist.functions}}

\begin{fulllineitems}
\phantomsection\label{\detokenize{index:hammindist.functions.max_clique_tomita}}
\pysigstartsignatures
\pysiglinewithargsret
{\sphinxcode{\sphinxupquote{hammindist.functions.}}\sphinxbfcode{\sphinxupquote{max\_clique\_tomita}}}
{\sphinxparam{\DUrole{n}{adj}}\sphinxparamcomma \sphinxparam{\DUrole{n}{N}}\sphinxparamcomma \sphinxparam{\DUrole{n}{vertices}}\sphinxparamcomma \sphinxparam{\DUrole{n}{use\_translation}\DUrole{o}{=}\DUrole{default_value}{True}}}
{}
\pysigstopsignatures
\sphinxAtStartPar
Tomita’s algorithm for the length of maximum clique
:param adj: Adjacency bitsets for the graph. adj{[}i{]} is an integer
\begin{quote}

\sphinxAtStartPar
where bit j indicates an edge between vertices i and j.
\end{quote}
\begin{quote}\begin{description}
\sphinxlineitem{Parameters}\begin{itemize}
\item {} 
\sphinxAtStartPar
\sphinxstyleliteralstrong{\sphinxupquote{N}} (\sphinxstyleliteralemphasis{\sphinxupquote{int}}) \textendash{} number of nodes in the graph

\item {} 
\sphinxAtStartPar
\sphinxstyleliteralstrong{\sphinxupquote{vertices}} (\sphinxstyleliteralemphasis{\sphinxupquote{List}}\sphinxstyleliteralemphasis{\sphinxupquote{{[}}}\sphinxstyleliteralemphasis{\sphinxupquote{any}}\sphinxstyleliteralemphasis{\sphinxupquote{{]}}}) \textendash{} nodes to work

\end{itemize}

\sphinxlineitem{Returns}
\sphinxAtStartPar
maximun size finded
best\_clique (List{[}int{]}): the best clique that the fuction was able to build

\sphinxlineitem{Return type}
\sphinxAtStartPar
max\_size (int)

\end{description}\end{quote}

\end{fulllineitems}

\index{max\_set() (in module hammindist.functions)@\spxentry{max\_set()}\spxextra{in module hammindist.functions}}

\begin{fulllineitems}
\phantomsection\label{\detokenize{index:hammindist.functions.max_set}}
\pysigstartsignatures
\pysiglinewithargsret
{\sphinxcode{\sphinxupquote{hammindist.functions.}}\sphinxbfcode{\sphinxupquote{max\_set}}}
{\sphinxparam{\DUrole{n}{codes\_list}}\sphinxparamcomma \sphinxparam{\DUrole{n}{min\_distance}}}
{}
\pysigstopsignatures
\sphinxAtStartPar
Build the set C wich contains all the strings that are a distance d between them
:param codes\_list: the strings, arrays or tuples that we want to know if are in the “best\_set”
:type codes\_list: List{[}any{]}
:param min\_distance: the distance that must have the elements of the “best\_set”
:type min\_distance: int
\begin{quote}\begin{description}
\sphinxlineitem{Returns}
\sphinxAtStartPar
the set of strings with min\_distance between them

\sphinxlineitem{Return type}
\sphinxAtStartPar
best\_set (set)

\end{description}\end{quote}

\end{fulllineitems}

\index{max\_set\_bounded() (in module hammindist.functions)@\spxentry{max\_set\_bounded()}\spxextra{in module hammindist.functions}}

\begin{fulllineitems}
\phantomsection\label{\detokenize{index:hammindist.functions.max_set_bounded}}
\pysigstartsignatures
\pysiglinewithargsret
{\sphinxcode{\sphinxupquote{hammindist.functions.}}\sphinxbfcode{\sphinxupquote{max\_set\_bounded}}}
{\sphinxparam{\DUrole{n}{length}}\sphinxparamcomma \sphinxparam{\DUrole{n}{min\_distance}}}
{}
\pysigstopsignatures
\sphinxAtStartPar
Build the set C wich contains all the strings that are a distance d between
them using a upper bound
:param codes\_list: the strings, arrays or tuples that we want to know if are in the “best\_set”
:type codes\_list: List{[}any{]}
:param min\_distance: the distance that must have the elements of the “best\_set”
:type min\_distance: int
\begin{quote}\begin{description}
\sphinxlineitem{Returns}
\sphinxAtStartPar
the set of strings with min\_distance between them

\sphinxlineitem{Return type}
\sphinxAtStartPar
best\_set (set)

\end{description}\end{quote}

\end{fulllineitems}

\index{plotkin\_bound() (in module hammindist.functions)@\spxentry{plotkin\_bound()}\spxextra{in module hammindist.functions}}

\begin{fulllineitems}
\phantomsection\label{\detokenize{index:hammindist.functions.plotkin_bound}}
\pysigstartsignatures
\pysiglinewithargsret
{\sphinxcode{\sphinxupquote{hammindist.functions.}}\sphinxbfcode{\sphinxupquote{plotkin\_bound}}}
{\sphinxparam{\DUrole{n}{length}}\sphinxparamcomma \sphinxparam{\DUrole{n}{distance}}}
{}
\pysigstopsignatures
\sphinxAtStartPar
Compute the Plotkin upper bound for A(n, d).
\begin{description}
\sphinxlineitem{Counts the total pairwise distance of a code in two ways:}\begin{itemize}
\item {} 
\sphinxAtStartPar
from below: each of the C(M,2) pairs contributes at least d,
so S \textgreater{}= M(M\sphinxhyphen{}1)/2 * d.

\item {} 
\sphinxAtStartPar
from above: each of the n bit positions contributes at most
M\textasciicircum{}2/4 differing pairs, so S \textless{}= n * M\textasciicircum{}2 / 4.

\end{itemize}

\end{description}

\sphinxAtStartPar
Combining both inequalities and solving for M yields
A(n, d) \textless{}= 2 * floor(d / (2d \sphinxhyphen{} n)), but only when 2d \textgreater{} n.
When 2d \textless{}= n the bound does not apply and infinity is returned.
\begin{quote}\begin{description}
\sphinxlineitem{Parameters}\begin{itemize}
\item {} 
\sphinxAtStartPar
\sphinxstyleliteralstrong{\sphinxupquote{length}} (\sphinxstyleliteralemphasis{\sphinxupquote{int}}) \textendash{} n, the length of the codes.

\item {} 
\sphinxAtStartPar
\sphinxstyleliteralstrong{\sphinxupquote{distance}} (\sphinxstyleliteralemphasis{\sphinxupquote{int}}) \textendash{} d, the minimum Hamming distance required.

\end{itemize}

\sphinxlineitem{Returns}
\sphinxAtStartPar
\begin{description}
\sphinxlineitem{the Plotkin upper bound on A(n, d),}
\sphinxAtStartPar
or float(‘inf’) if the bound does not apply.

\end{description}


\sphinxlineitem{Return type}
\sphinxAtStartPar
int | float

\end{description}\end{quote}

\end{fulllineitems}

\index{select\_words\_to\_remove() (in module hammindist.functions)@\spxentry{select\_words\_to\_remove()}\spxextra{in module hammindist.functions}}

\begin{fulllineitems}
\phantomsection\label{\detokenize{index:hammindist.functions.select_words_to_remove}}
\pysigstartsignatures
\pysiglinewithargsret
{\sphinxcode{\sphinxupquote{hammindist.functions.}}\sphinxbfcode{\sphinxupquote{select\_words\_to\_remove}}}
{\sphinxparam{\DUrole{n}{current\_code}}\sphinxparamcomma \sphinxparam{\DUrole{n}{tabu\_expiry}}\sphinxparamcomma \sphinxparam{\DUrole{n}{step}}\sphinxparamcomma \sphinxparam{\DUrole{n}{k}}}
{}
\pysigstopsignatures
\sphinxAtStartPar
Select k codewords to remove.

\sphinxAtStartPar
Priority is given to codewords that are not currently tabu.
If there are not enough non\sphinxhyphen{}tabu codewords, the remaining
selections are completed with tabu codewords.
:param current\_code: List of codewords currently in the code.
\begin{quote}

\sphinxAtStartPar
Each codeword is typically an integer
representing a binary string.
\end{quote}
\begin{quote}\begin{description}
\sphinxlineitem{Parameters}\begin{itemize}
\item {} 
\sphinxAtStartPar
\sphinxstyleliteralstrong{\sphinxupquote{tabu\_expiry}} (\sphinxstyleliteralemphasis{\sphinxupquote{Dict}}\sphinxstyleliteralemphasis{\sphinxupquote{{[}}}\sphinxstyleliteralemphasis{\sphinxupquote{Any}}\sphinxstyleliteralemphasis{\sphinxupquote{, }}\sphinxstyleliteralemphasis{\sphinxupquote{int}}\sphinxstyleliteralemphasis{\sphinxupquote{{]}}}) \textendash{} Dictionary mapping codewords to the
step number when their tabu status expires.
A codeword is considered tabu if its expiry
step is greater than the current step.

\item {} 
\sphinxAtStartPar
\sphinxstyleliteralstrong{\sphinxupquote{step}} (\sphinxstyleliteralemphasis{\sphinxupquote{int}}) \textendash{} Current iteration step number. Used to determine which
codewords have expired tabu status.

\item {} 
\sphinxAtStartPar
\sphinxstyleliteralstrong{\sphinxupquote{k}} (\sphinxstyleliteralemphasis{\sphinxupquote{int}}) \textendash{} Number of codewords to select for removal.

\end{itemize}

\sphinxlineitem{Returns}
\sphinxAtStartPar
\begin{description}
\sphinxlineitem{A list of k codewords selected for removal. The selection}
\sphinxAtStartPar
prioritizes non\sphinxhyphen{}tabu codewords, but may include tabu
codewords if necessary to reach k selections.

\end{description}


\sphinxlineitem{Return type}
\sphinxAtStartPar
List{[}Any{]}

\end{description}\end{quote}

\end{fulllineitems}

\index{simulated\_annealing() (in module hammindist.functions)@\spxentry{simulated\_annealing()}\spxextra{in module hammindist.functions}}

\begin{fulllineitems}
\phantomsection\label{\detokenize{index:hammindist.functions.simulated_annealing}}
\pysigstartsignatures
\pysiglinewithargsret
{\sphinxcode{\sphinxupquote{hammindist.functions.}}\sphinxbfcode{\sphinxupquote{simulated\_annealing}}}
{\sphinxparam{\DUrole{n}{n}}\sphinxparamcomma \sphinxparam{\DUrole{n}{d}}\sphinxparamcomma \sphinxparam{\DUrole{n}{iterations}\DUrole{o}{=}\DUrole{default_value}{5000}}\sphinxparamcomma \sphinxparam{\DUrole{n}{T\_start}\DUrole{o}{=}\DUrole{default_value}{1.0}}\sphinxparamcomma \sphinxparam{\DUrole{n}{T\_end}\DUrole{o}{=}\DUrole{default_value}{0.01}}\sphinxparamcomma \sphinxparam{\DUrole{n}{tabu\_tenure}\DUrole{o}{=}\DUrole{default_value}{20}}\sphinxparamcomma \sphinxparam{\DUrole{n}{max\_perturbation}\DUrole{o}{=}\DUrole{default_value}{3}}\sphinxparamcomma \sphinxparam{\DUrole{n}{seed}\DUrole{o}{=}\DUrole{default_value}{None}}}
{}
\pysigstopsignatures
\sphinxAtStartPar
Estimate A(n, d) using simulated annealing with tabu memory.

\sphinxAtStartPar
At each iteration:
1. Remove between 1 and max\_perturbation codewords.
2. Mark removed codewords as tabu for a fixed tenure.
3. Attempt to add up to max\_perturbation + 1 new codewords.
4. Accept or reject the candidate solution according to
\begin{quote}

\sphinxAtStartPar
the simulated annealing criterion.
\end{quote}

\sphinxAtStartPar
A reheating mechanism is also applied whenever no improvement
is observed for an extended number of iterations.

\end{fulllineitems}

\index{sphere\_packing\_bound() (in module hammindist.functions)@\spxentry{sphere\_packing\_bound()}\spxextra{in module hammindist.functions}}

\begin{fulllineitems}
\phantomsection\label{\detokenize{index:hammindist.functions.sphere_packing_bound}}
\pysigstartsignatures
\pysiglinewithargsret
{\sphinxcode{\sphinxupquote{hammindist.functions.}}\sphinxbfcode{\sphinxupquote{sphere\_packing\_bound}}}
{\sphinxparam{\DUrole{n}{length}}\sphinxparamcomma \sphinxparam{\DUrole{n}{distance}}}
{}
\pysigstopsignatures
\sphinxAtStartPar
Compute the sphere\sphinxhyphen{}packing (Hamming) upper bound for A(n, d).

\sphinxAtStartPar
Treats each codeword as the center of a Hamming ball of radius
t = (d\sphinxhyphen{}1)//2. Because balls around distinct codewords are
disjoint, their total volume cannot exceed the full space 2\textasciicircum{}n,
giving A(n, d) \textless{}= 2\textasciicircum{}n / sum\_\{i=0\}\textasciicircum{}\{t\} C(n, i).
\begin{quote}\begin{description}
\sphinxlineitem{Parameters}\begin{itemize}
\item {} 
\sphinxAtStartPar
\sphinxstyleliteralstrong{\sphinxupquote{length}} (\sphinxstyleliteralemphasis{\sphinxupquote{int}}) \textendash{} n, the length of the codes.

\item {} 
\sphinxAtStartPar
\sphinxstyleliteralstrong{\sphinxupquote{distance}} (\sphinxstyleliteralemphasis{\sphinxupquote{int}}) \textendash{} d, the minimum Hamming distance required.

\end{itemize}

\sphinxlineitem{Returns}
\sphinxAtStartPar
the sphere\sphinxhyphen{}packing upper bound on A(n, d).

\sphinxlineitem{Return type}
\sphinxAtStartPar
int

\end{description}\end{quote}

\end{fulllineitems}

\index{try\_add\_codewords() (in module hammindist.functions)@\spxentry{try\_add\_codewords()}\spxextra{in module hammindist.functions}}

\begin{fulllineitems}
\phantomsection\label{\detokenize{index:hammindist.functions.try_add_codewords}}
\pysigstartsignatures
\pysiglinewithargsret
{\sphinxcode{\sphinxupquote{hammindist.functions.}}\sphinxbfcode{\sphinxupquote{try\_add\_codewords}}}
{\sphinxparam{\DUrole{n}{candidate}}\sphinxparamcomma \sphinxparam{\DUrole{n}{pool}}\sphinxparamcomma \sphinxparam{\DUrole{n}{d}}\sphinxparamcomma \sphinxparam{\DUrole{n}{k}}}
{}
\pysigstopsignatures
\sphinxAtStartPar
Attempt to add up to k + 1 compatible codewords.

\sphinxAtStartPar
A codeword is added only if it preserves the required
minimum distance from all codewords already in the solution.
:param candidate: Current solution list of codewords already
\begin{quote}

\sphinxAtStartPar
selected. This list will be modified in\sphinxhyphen{}place.
\end{quote}
\begin{quote}\begin{description}
\sphinxlineitem{Parameters}\begin{itemize}
\item {} 
\sphinxAtStartPar
\sphinxstyleliteralstrong{\sphinxupquote{pool}} (\sphinxstyleliteralemphasis{\sphinxupquote{List}}\sphinxstyleliteralemphasis{\sphinxupquote{{[}}}\sphinxstyleliteralemphasis{\sphinxupquote{int}}\sphinxstyleliteralemphasis{\sphinxupquote{{]}}}) \textendash{} List of available codewords to consider adding.
These are typically from a neighborhood or
perturbation operation.

\item {} 
\sphinxAtStartPar
\sphinxstyleliteralstrong{\sphinxupquote{d}} (\sphinxstyleliteralemphasis{\sphinxupquote{int}}) \textendash{} Required minimum Hamming distance. Any new codeword must
have distance at least d from all codewords already in
the candidate solution.

\item {} 
\sphinxAtStartPar
\sphinxstyleliteralstrong{\sphinxupquote{k}} (\sphinxstyleliteralemphasis{\sphinxupquote{int}}) \textendash{} Maximum number of codewords to add beyond the first one.
The function will add at most k+1 codewords total.

\end{itemize}

\sphinxlineitem{Returns}
\sphinxAtStartPar
\begin{description}
\sphinxlineitem{The modified candidate solution list (same list object}
\sphinxAtStartPar
as the input, modified in\sphinxhyphen{}place) with newly added
compatible codewords appended.

\end{description}


\sphinxlineitem{Return type}
\sphinxAtStartPar
List{[}int{]}

\end{description}\end{quote}

\end{fulllineitems}

\index{upper\_bound() (in module hammindist.functions)@\spxentry{upper\_bound()}\spxextra{in module hammindist.functions}}

\begin{fulllineitems}
\phantomsection\label{\detokenize{index:hammindist.functions.upper_bound}}
\pysigstartsignatures
\pysiglinewithargsret
{\sphinxcode{\sphinxupquote{hammindist.functions.}}\sphinxbfcode{\sphinxupquote{upper\_bound}}}
{\sphinxparam{\DUrole{n}{length}}\sphinxparamcomma \sphinxparam{\DUrole{n}{distance}}}
{}
\pysigstopsignatures
\sphinxAtStartPar
Compute the tightest available upper bound for A(n, d).

\sphinxAtStartPar
Takes the minimum of the sphere\sphinxhyphen{}packing (Hamming) bound and the
Plotkin bound. When Plotkin does not apply (2d \textless{}= n) its value is
infinity, so the sphere\sphinxhyphen{}packing bound is returned automatically.
\begin{quote}\begin{description}
\sphinxlineitem{Parameters}\begin{itemize}
\item {} 
\sphinxAtStartPar
\sphinxstyleliteralstrong{\sphinxupquote{length}} (\sphinxstyleliteralemphasis{\sphinxupquote{int}}) \textendash{} n, the word length in bits.

\item {} 
\sphinxAtStartPar
\sphinxstyleliteralstrong{\sphinxupquote{distance}} (\sphinxstyleliteralemphasis{\sphinxupquote{int}}) \textendash{} d, the minimum Hamming distance required.

\end{itemize}

\sphinxlineitem{Returns}
\sphinxAtStartPar
the tightest upper bound on A(n, d) from the two methods.

\sphinxlineitem{Return type}
\sphinxAtStartPar
int

\end{description}\end{quote}

\end{fulllineitems}



\renewcommand{\indexname}{Python Module Index}
\begin{sphinxtheindex}
\let\bigletter\sphinxstyleindexlettergroup
\bigletter{h}
\item\relax\sphinxstyleindexentry{hammindist.functions}\sphinxstyleindexpageref{index:\detokenize{module-hammindist.functions}}
\end{sphinxtheindex}

\renewcommand{\indexname}{Index}
\printindex
\end{document}